\hypertarget{data-custody-preserves-the-meaning-of-a-causal-claim}{%
\section{Data custody preserves the meaning of a causal
claim}\label{data-custody-preserves-the-meaning-of-a-causal-claim}}

\label{exec14:artifact-custody}

\hypertarget{the-object-to-preserve-is-a-decision-with-its-information}{%
\subsection{The object to preserve is a decision with its
information}\label{the-object-to-preserve-is-a-decision-with-its-information}}

A servicing analysis can be arithmetically reproducible and still answer
the wrong question. Suppose an analyst computes the effect of an offer
on the next twelve months of account value among accounts eligible at
the decision date. A second workflow receives the same numerical moments
but reports a twenty-four-month effect for all customers. No file
corruption is needed for that mistake. Population, outcome horizon,
intervention, information clock, units and sampling law belong to the
scientific object being transferred. They must travel with the data,
graph and result. A consumer that cannot establish this agreement should
return an unavailable result or request an explicit new analysis.

Let \(C\) denote the declared scientific context, including those fields
and the graph with its assumption statuses. Let \(S\) be the local
deterministic JSON serialization, \(H\) the SHA-256 digest, and
\(B=S(V)\) the stored bytes for a JSON value \(V\). The implemented
local store uses a descriptor

\[
D=(H(B),\;\text{media type},\;\text{schema},\;\text{role},\;|B|,\;\text{custody}),
\]

and derives a handle from \(H(S(D))\), preceded by a fixed letter. A
scientific input also carries \(H(S(C))\) inside its payload. These are
separate bindings: the first identifies bytes and their declared
interpretation, while the second ties a payload to a particular
scientific request. The store's serialization sorts object keys and
excludes nonfinite JSON numbers. It is an inspected Python serialization
rule, not an established canonicalization standard for every language or
a proof that distinct textual descriptions express the same causal
question.~\citep{exec14_custody_implementation}

For instance, a moment summary contains both its numbers and the context
digest. Changing the horizon from twelve to twenty-four months changes
the context digest even when all moments are unchanged. The service
checks the expected artifact role and schema, reads and verifies the
referenced bytes, and compares the embedded context digest. In an actual
local trial with a generated quadratic sensitivity input, the
twelve-month request computed an empty sublevel set; the otherwise
identical twenty-four-month request raised a context error. Through the
subprocess runtime the latter became a failed, unavailable response with
no numerical result. This distinction matters: an empty set is a
mathematical outcome for a supplied problem; an unavailable result says
that the requested problem was not computed. Conflating them would turn
a custody failure into apparent scientific evidence.
\citep{exec14_custody_observations}

Hash agreement does not establish that the population declaration is
true, that timestamps reflect actual availability, or that a supplied
graph is causal. The declaration and the authority behind it require
different evidence. A data owner can provide a signed account of
extraction and eligibility, but the signature would establish who made
that account, not the validity of ignorability, exclusion or a
measurement model. The proposed product should preserve these
distinctions from ingestion through the reported monetary consequence.

\hypertarget{what-the-local-implementation-establishes}{%
\subsection{What the local implementation
establishes}\label{what-the-local-implementation-establishes}}

The reference store is restricted to nonsymlink runtime paths inside
this repository. It stores finite JSON values up to eight MiB and
publishes each content or metadata file atomically, with a directory
flush. Reading checks agreement between the supplied reference and
recorded metadata, file length and the byte digest. The operation
dispatcher adds role, schema and declared scientific-context checks. The
runtime verifies input references before execution and before
considering an idempotent replay. \citep{exec14_custody_implementation}

Fifteen generated local cases exercise this boundary. Repeated insertion
under the same role returned the same reference; changing only the role
retained the content digest and changed the handle. Changed bytes,
changed metadata, an incompatible horizon, a wrong role or schema, an
insufficient read limit, a nonfinite number, and selected symlink paths
were refused. The two existing service/runtime test modules also passed
all sixteen tests in this environment. These observations support
selected implementation claims; they do not estimate an operational
error rate, prove correctness against every possible input, or validate
bank records. \citep{exec14_custody_observations}

Replay deserves its own scientific interpretation. A new request
correlation identifier may retrieve the same result, but reusing the
idempotency key for a changed horizon returned a conflict. After the
referenced input was altered, replay returned an artifact-access failure
with no result. Its key binds the request digest, selected
implementation and schema hashes, Python version and declared dependency
versions. A current source digest prevents an ordinary unchanged-key
replay across a detected implementation change; it is not a complete
software supply-chain attestation. The local operator still controls
these files and the environment.
\citep{exec14_custody_implementation,exec14_custody_observations}

The report path has an equally narrow role. Its summaries explicitly
carry a schema-only verification label. Preserving a canonical response
in a report does not promote execution evidence into accepted causal
inference. For example, an assumption-conditional identification
expression must retain that qualification even if its artifact passes
every byte check. Similarly, a refused or unsupported operation must not
disappear from a report that counts only successful tasks. The preserved
runtime tests establish selected report and cancellation behavior;
institutional report approval remains an additional
decision.~\citep{exec14_custody_implementation,exec14_custody_observations}

\hypertarget{two-observed-boundaries-and-the-repairs-they-require}{%
\subsection{Two observed boundaries and the repairs they
require}\label{two-observed-boundaries-and-the-repairs-they-require}}

The first boundary concerns the relation between a handle and its
descriptor. The writer derives the handle from the descriptor, but the
current reader does not independently recompute that derivation. A
controlled local trial changed a generated balance from 100 to 900,
rewrote the stored descriptor with the new byte digest and length, and
supplied the correspondingly changed reference under the original
handle. The old trusted reference was refused. The new supplied
reference was accepted, although recomputing the descriptor digest
produced a different handle. The test required write access to content
and metadata and control of the supplied reference. It demonstrates a
limit of the single-operator model; it does not establish an exposed
remote attack path.~\citep{exec14_custody_observations}

There is a direct local repair: validate the complete descriptor schema
and recompute its handle before returning content or accepting stored
metadata. Require the derived handle to equal both the requested
identifier and the identifier in the stored reference. Preserve the
original implementation as the evidence baseline, and qualify a
successor against the old trusted reference, a coherently rewritten
reference, a mismatched metadata identity, missing descriptor fields and
extra fields. This is a proposed repair, not an implemented feature. It
also does not solve unrestricted substitution: a writer who can replace
the trusted request can supply an entirely new, internally consistent
artifact and handle. Preventing that substitution requires an
independently trusted manifest or authenticated authorization binding
the requester, purpose and allowed inputs. \citep{exec14_custody_design}

The second boundary concerns availability. The reader opens a path
before checking that its file descriptor denotes a bounded regular file.
The open uses a flag rejecting a final symlink, but lacks nonblocking
mode. In a disposable subprocess, opening a named pipe without a writer
did not return within a one-second deadline after the child signalled
that it was about to read. The subprocess was terminated. The observed
block occurred before the regular-file test could reject the pipe.
Because the runtime performs its initial reference verification in the
parent, a deadline on the later worker does not by itself bound that
initial access. This is an actual local availability finding, with no
claim about deployment reachability or a measured service-level failure
rate. \citep{exec14_custody_implementation,exec14_custody_observations}

A candidate repair opens candidate paths without following a final
symlink and in nonblocking mode, immediately checks the returned
descriptor, and accepts only bounded regular files. The acceptance cases
must include pipes with and without writers, directories, symlinks and
oversized or changing files; each rejection must occur before any causal
worker starts. Directory file descriptors and controlled directory
ownership should also be reviewed if the threat model includes
concurrent path replacement. A prior path resolution check alone does
not establish race-free access. Nonblocking open is a concrete first
repair for the demonstrated pipe, not a promise to bound all filesystem
delays or solve every race. A production design still needs an access
deadline and an operating-system-specific review.
\citep{exec14_custody_design}

These findings do not invalidate the observed ordinary mutation
refusals. They define the assumptions under which those refusals are
useful. For a single researcher working with generated inputs, agreement
with a trusted reference can detect accidental changes. For multiple
operators, automated agents or a bank repository, the trust anchor and
authorization decision must live beyond a request's freely supplied
custody string.

\hypertarget{custody-across-languages-repositories-and-agent-tools}{%
\subsection{Custody across languages, repositories and agent
tools}\label{custody-across-languages-repositories-and-agent-tools}}

The transfer contract should distinguish a logical artifact from a path
on one machine. A Python caller, an R adapter and a Java discovery
engine may share a causal task without sharing a filesystem. The
prospective adapter resolves an allowed handle within its own custody
boundary, checks the manifest and exact schema, and obtains the material
needed for its declared operation. It preserves graph endpoints,
effect-set topology, units, population, clock and uncertainty semantics
in the returned object. Copying a path into a CLI flag or an MCP request
does not establish that agreement. No new sibling-repository execution
is claimed by this design. \citep{exec14_custody_design}

Two prospective arrangements have different costs. A controlled local
service may expose a small allowlist of typed handles and operations to
trusted callers. This keeps initial experiments inexpensive but depends
on the service's operator and host boundary. A broader artifact service
would bind authenticated principals and permitted purposes to manifests,
retain authorized provenance, and enforce current access and revocation
decisions. That arrangement adds operational complexity and requires an
owner for identity, keys, access policies and lifecycle events. The
product proposal can specify and compare them before claiming either is
an accepted bank service.

A context digest provides a strict equality check, which can be
desirable at a computational boundary but does not resolve semantic
equivalence. Equivalent population wording or a different serialization
of the same declared numerical quantity may fail equality; different
underlying data can also carry an identical but false declaration. A
future interoperable context schema therefore needs a versioned
representation, explicit units and clock conventions, and reviewed
migrations. If an adapter changes a horizon, aggregation or currency
conversion, it should produce a new artifact with the transformation and
input identities recorded. It should not overwrite the meaning attached
to an existing result.

Retention creates another distinction between repeatability and
permission. An immutable result can remain scientifically reproducible
while its source records become unavailable for a later request. The
service should check current access policy before replay, distinguish
revoked access from missing bytes, and retain only the audit material
its owner permits. An invalidated source should have an explicit status
and downstream impact record. Deletion, backups, redaction, derived
artifacts and legally permitted audit retention require owner decisions.
The current local store has no qualified revocation or institutional
retention service. Designing one requires more than continuing to
compare hashes.~\citep{exec14_custody_design}

\hypertarget{an-affordable-qualification-path}{%
\subsection{An affordable qualification
path}\label{an-affordable-qualification-path}}

The next experiment can use generated accounts, public documents and
deliberately altered copies. It should first qualify the two
demonstrated local repairs without changing the original evidence. Then
it can exercise an adapter against distinct populations, horizons,
endpoint types and information clocks, including refusals and missing
inputs. Public data can test acquisition and transformation lineage, but
a public label does not automatically establish redistribution rights or
causal ground truth. Each chosen dataset needs its own provenance and
permitted-use review.

The prospective acceptance cases divide into concrete responsibilities:

\begin{longtable}[]{@{}lll@{}}
\toprule
\begin{minipage}[b]{0.30\columnwidth}\raggedright
Requirement\strut
\end{minipage} & \begin{minipage}[b]{0.30\columnwidth}\raggedright
Local or public-evidence next step\strut
\end{minipage} & \begin{minipage}[b]{0.30\columnwidth}\raggedright
Additional owner input\strut
\end{minipage}\tabularnewline
\midrule
\endhead
\begin{minipage}[t]{0.30\columnwidth}\raggedright
Byte, descriptor and handle consistency\strut
\end{minipage} & \begin{minipage}[t]{0.30\columnwidth}\raggedright
Qualify a successor reader against ordinary and coherent replacement
cases\strut
\end{minipage} & \begin{minipage}[t]{0.30\columnwidth}\raggedright
Trusted manifest location and who may publish it\strut
\end{minipage}\tabularnewline
\begin{minipage}[t]{0.30\columnwidth}\raggedright
Bounded input access\strut
\end{minipage} & \begin{minipage}[t]{0.30\columnwidth}\raggedright
Reject unexpected file types before worker launch; test deadlines and
controlled path replacement\strut
\end{minipage} & \begin{minipage}[t]{0.30\columnwidth}\raggedright
Target filesystem, latency limits and operating model\strut
\end{minipage}\tabularnewline
\begin{minipage}[t]{0.30\columnwidth}\raggedright
Scientific meaning\strut
\end{minipage} & \begin{minipage}[t]{0.30\columnwidth}\raggedright
Preserve role, schema, context, endpoint and uncertainty contracts in
adapter fixtures\strut
\end{minipage} & \begin{minipage}[t]{0.30\columnwidth}\raggedright
Population, time, assignment and measurement declarations with
accountable provenance\strut
\end{minipage}\tabularnewline
\begin{minipage}[t]{0.30\columnwidth}\raggedright
Authorization and agent scope\strut
\end{minipage} & \begin{minipage}[t]{0.30\columnwidth}\raggedright
Model permitted and forbidden handles, purposes and operations using
generated principals\strut
\end{minipage} & \begin{minipage}[t]{0.30\columnwidth}\raggedright
Identity authority, policy owner, key custody and access rules\strut
\end{minipage}\tabularnewline
\begin{minipage}[t]{0.30\columnwidth}\raggedright
Retention and revocation\strut
\end{minipage} & \begin{minipage}[t]{0.30\columnwidth}\raggedright
Exercise expired and revoked input scenarios without silently replaying
results\strut
\end{minipage} & \begin{minipage}[t]{0.30\columnwidth}\raggedright
Retention, deletion, backups and permitted audit evidence\strut
\end{minipage}\tabularnewline
\begin{minipage}[t]{0.30\columnwidth}\raggedright
Cross-repository use\strut
\end{minipage} & \begin{minipage}[t]{0.30\columnwidth}\raggedright
Compare actual local behavior with the declared adapter contract\strut
\end{minipage} & \begin{minipage}[t]{0.30\columnwidth}\raggedright
Named consumer, supported environment and acceptance decision\strut
\end{minipage}\tabularnewline
\bottomrule
\end{longtable}

The bank need not fund an authentic institutional trial to establish the
value of these checks. A prototype can first show that a wrong horizon
is refused, an unavailable result stays unavailable, and failed access
cannot masquerade as an empty effect set. Its success criterion is
preserved meaning and a documented disposition for every request. Actual
bank provenance, access authority, operating performance and adoption
remain later claims requiring their own
evidence.~\citep{exec14_custody_design}

\label{exec14:artifact-custody-end}
